\documentclass[11pt]{article}
\usepackage[margin=1.85cm]{geometry}
\usepackage{amsmath,amssymb}
\usepackage{booktabs}
\usepackage{xcolor}
\usepackage{listings}
\usepackage{array}
\usepackage{enumitem}
\usepackage[hidelinks]{hyperref}
\usepackage{titlesec}

\definecolor{codebg}{RGB}{246,247,249}
\definecolor{codekw}{RGB}{0,88,158}
\definecolor{codecm}{RGB}{105,115,105}
\definecolor{codestr}{RGB}{158,60,20}
\definecolor{rulegray}{RGB}{205,210,215}

\lstdefinestyle{ml}{
  language=Matlab,
  backgroundcolor=\color{codebg},
  basicstyle=\ttfamily\footnotesize,
  keywordstyle=\color{codekw}\bfseries,
  commentstyle=\color{codecm}\itshape,
  stringstyle=\color{codestr},
  showstringspaces=false,
  breaklines=true,
  breakatwhitespace=false,
  columns=fullflexible,
  keepspaces=true,
  frame=single,
  framesep=4pt,
  rulecolor=\color{rulegray},
  aboveskip=6pt,
  belowskip=6pt
}
\lstdefinestyle{inl}{
  language=Matlab,
  basicstyle=\ttfamily\small,
  keywordstyle=\color{codekw},
  commentstyle=\color{codecm}\itshape,
  stringstyle=\color{codestr},
  showstringspaces=false,
  frame=none,
  backgroundcolor=\color{white},
  columns=fullflexible,
  keepspaces=true
}
\lstset{style=ml}

\titleformat{\section}{\large\bfseries\color{codekw}}{\thesection}{0.5em}{}
\titlespacing*{\section}{0pt}{10pt}{4pt}
\titleformat{\subsection}{\normalsize\bfseries}{\thesubsection}{0.5em}{}
\titlespacing*{\subsection}{0pt}{8pt}{3pt}

\setlength{\parindent}{0pt}
\setlength{\parskip}{4pt}

\begin{document}

\begin{center}
{\LARGE\bfseries MATLAB cheat sheet}\\[2pt]
{\large ECON*6020 --- Assignment 1: Canadian business cycle data, 1976Q1--2026Q2}\\[4pt]
{\small Data file \texttt{Assignment1\_data.xlsx}, sheet \texttt{Combined}: 202 quarters $\times$ 66 columns.\\
Companion to \texttt{assignment1.m}. This sheet lists the \emph{syntax} you need --- not the answers.}
\end{center}

\vspace{2pt}
\hrule
\vspace{6pt}

\section{Load the file and look at it}

\begin{lstlisting}
clear; clc; close all

T = readtable("Assignment1_data.xlsx", Sheet = "Combined", ...
              VariableNamingRule = "preserve");

disp(head(T, 8))                            % first 8 rows of the sheet
fprintf("%d quarters x %d columns\n", height(T), width(T));
var = string(T.Properties.VariableNames);    % all 66 names as a string vector
q   = string(T.quarter);                     % "1976Q1" ... "2026Q2"  (column 1)
Tq  = height(T);                             % 202
\end{lstlisting}

\lstinline[style=inl]!readtable! reads the sheet into a \emph{table}; a column is then
\lstinline[style=inl]!T.name!. Things worth knowing: \lstinline[style=inl]!height!/\lstinline[style=inl]!width! give rows/columns
(\lstinline[style=inl]!size! works too), \lstinline[style=inl]!head(T,n)! shows the first $n$ rows, \lstinline[style=inl]!T{:,k}!
gives column $k$ as a plain numeric vector, and \lstinline[style=inl]!string(...)! converts to
text so \lstinline[style=inl]!==! comparisons work.

\section{How the column names are built}

Every GDP column name is three segments joined by \lstinline[style=inl]!_!:

\begin{center}
\lstinline[style=inl]!basis! \ \lstinline[style=inl]!_! \ \lstinline[style=inl]!seasonally_adjusted_at_annual_rates! \ \lstinline[style=inl]!_! \ \lstinline[style=inl]!what!
\end{center}

where \lstinline[style=inl]!basis! is \lstinline[style=inl]!current_prices! (nominal) or
\lstinline[style=inl]!chained_2017_dollars! (real), and \lstinline[style=inl]!what! is one of the 31 components:

\vspace{-2pt}
{\footnotesize
\begin{center}
\begin{tabular}{@{}ll@{}}
final\_consumption\_expenditure & exports\_of\_goods\_and\_services \\
household\_final\_consumption\_expenditure & exports\_of\_goods \\
goods & exports\_of\_services \\
durable\_goods & less\_imports\_of\_goods\_and\_services \\
semi\_durable\_goods & imports\_of\_goods \\
non\_durable\_goods & imports\_of\_services \\
services & statistical\_discrepancy \\
non\_profit\_institutions\_serving\_households\_final\_consumption\_expenditure & gross\_domestic\_product\_at\_market\_prices \\
general\_governments\_final\_consumption\_expenditure & final\_domestic\_demand \\
gross\_fixed\_capital\_formation & investment\_in\_inventories \\
business\_gross\_fixed\_capital\_formation & of\_which\_business\_investment\_in\_inventories \\
residential\_structures & non\_farm \\
non\_residential\_structures\_machinery\_and\_equipment & farm \\
non\_residential\_structures & \\
machinery\_and\_equipment & \\
intellectual\_property\_products & \\
non\_profit\_institutions\_serving\_households\_gross\_fixed\_capital\_formation & \\
general\_governments\_gross\_fixed\_capital\_formation & \\
\end{tabular}
\end{center}}

The three extra columns are
\lstinline[style=inl]!hours_worked_v4391505_quarterly_sum_hours!,
\lstinline[style=inl]!population_v1_persons! and
\lstinline[style=inl]!exchange_rate_v37426_to_2016_v111666275_from_2017_quarterly_average_cad_per_usd!.

\subsection*{Build a name instead of typing it}

\begin{lstlisting}
nm = @(basis, what) join([basis, "seasonally_adjusted_at_annual_rates", what], "_");
R  = "chained_2017_dollars";     % real, 2017 dollars
N  = "current_prices";           % nominal

y    = T.chained_2017_dollars_seasonally_adjusted_at_annual_rates_gross_domestic_product_at_market_prices;
ynom = T.current_prices_seasonally_adjusted_at_annual_rates_gross_domestic_product_at_market_prices;
                                                                  % dot access works...
y    = T.(nm(R, "gross_domestic_product_at_market_prices"));      % ...or computed
ynom = T.(nm(N, "gross_domestic_product_at_market_prices"));
\end{lstlisting}

\subsection*{Every name is a legal MATLAB identifier}

\lstinline[style=inl]!isvarname! returns \lstinline[style=inl]!1! for all 66 names (checked), so all three
ways of reaching a column work:

\begin{lstlisting}
pop = T.population_v1_persons;                        % plain dot access
pop = T.("population_v1_persons");                    % same column, brackets
pop = T.(nm("chained_2017_dollars", "services"));     % computed from the pieces
\end{lstlisting}

The GDP names run up to 129 characters, so the computed form is usually less
typing; when you do type one out, copy it from \lstinline[style=inl]!T.Properties.VariableNames!
or from the header row of the Excel file.

\section{The 14 variables of the table}

\begin{lstlisting}
% --- 10 series straight from the file -------------------------------------
y    = T.(nm(R, "gross_domestic_product_at_market_prices"));      % output, real
c    = T.(nm(R, "semi_durable_goods")) + ...
       T.(nm(R, "non_durable_goods"))  + T.(nm(R, "services"));   % consumption
i    = T.(nm(R, "non_residential_structures_machinery_and_equipment"));
gc   = T.(nm(R, "general_governments_final_consumption_expenditure"));
gi   = T.(nm(R, "general_governments_gross_fixed_capital_formation"));
ex   = T.(nm(R, "exports_of_goods_and_services"));
im   = T.(nm(R, "less_imports_of_goods_and_services"));           % positive
n    = T.hours_worked_v4391505_quarterly_sum_hours;               % hours
pop  = T.population_v1_persons;                                   % persons
fx   = 100 * T.exchange_rate_v37426_to_2016_v111666275_from_2017_quarterly_average_cad_per_usd;
                                                                  % CAD/USD -> cents

% --- 4 variables you build ------------------------------------------------
nx   = 1 + (ex - im) ./ y;                  % net exports, 1 + NX/Y
w    = y ./ n;                              % labour productivity, Y/N
pgdp = 100 * ynom ./ y;                     % price = GDP deflator

delta = 0.025;                              % quarterly depreciation
gI    = mean(diff(i) ./ i(1:end-1));        % avg growth of investment
k     = zeros(Tq, 1);
k(1)  = i(1) / (gI + delta);                % K0 from (g+delta)K = I
for t = 1:Tq-1
    k(t+1) = i(t) + (1 - delta) * k(t);     % law of motion
end
mpk  = y ./ k;                              % capital productivity, Y/K

% --- one matrix, table order, then per capita ------------------------------
X = [y c i gc gi ex im nx n w k mpk fx pgdp];      % 202 x 14
QC = [1 2 3 4 5 6 7 9 11];                         % y c i gc gi ex im n k
X(:, QC) = X(:, QC) ./ pop;    % ratios (nx w mpk) and indexes (fx pgdp) are
                               % already per-capita, so leave them alone
\end{lstlisting}

Order of \lstinline[style=inl]!X! (this is the table's order): 1 output, 2 consumption,
3 investment, 4 gov consumption, 5 gov investment, 6 exports, 7 imports,
8 net exports, 9 hours, 10 labour productivity, 11 capital, 12 capital
productivity, 13 nominal exchange rate, 14 price.

\section{Operators: matrix algebra vs element-by-element}

\begin{center}
\begin{tabular}{@{}lll@{}}
\toprule
 & matrix algebra & element-by-element \\
\midrule
multiply & \texttt{A*B} & \texttt{A.*B} \\
divide   & \texttt{A/B} (solve) & \texttt{A./B} \\
power    & \texttt{A\^{}2} & \texttt{A.\^{}2} \\
\bottomrule
\end{tabular}
\end{center}

Rule of thumb: two series of the same length are \emph{always} \lstinline[style=inl]!./! and
\lstinline[style=inl]!.*!. The plain versions are for matrix products, and using them by mistake
gives either an error (\emph{Incorrect dimensions}) or a silent scalar matrix
product.

\section{HP filter (the method in the course's \texttt{hpfilter.m})}

\begin{lstlisting}
raw    = log(X);                                  % log FIRST
lambda = 1600;
d = (1+6*lambda)*ones(Tq,1);  x = diag(d);
d = -4*lambda*ones(Tq-1,1);   x = x + diag(d,1) + diag(d,-1);
d =  lambda*ones(Tq-2,1);     x = x + diag(d,2) + diag(d,-2);
x(1,1) = 1+lambda;      x(1,2) = -2*lambda;   x(2,1) = -2*lambda;
x(2,2) = 1+5*lambda;    x(Tq,Tq) = 1+lambda;  x(Tq-1,Tq) = -2*lambda;
x(Tq,Tq-1) = -2*lambda; x(Tq-1,Tq-1) = 1+5*lambda;
hpmat  = inv(x);                                  % HP matrix, built once
trend  = hpmat * raw;                             % trend, in logs
dev    = raw - trend;                             % cyclical component
\end{lstlisting}

The course file reads \lstinline[style=inl]!data2026q2.csv! and derives the number of columns
itself, so to reuse it exactly:
\lstinline[style=inl]!writematrix(X, "data2026q2.csv")! then run \lstinline[style=inl]!hpfilter!. It hard-codes
$T=202$, which is our sample length, so the sizes work out.

\section{Statistics}

\begin{lstlisting}
vcv  = (dev' * dev) / Tq;               % variance-covariance matrix (14x14)
v    = sum(vcv .* eye(14));             % variances on the diagonal
sd   = sqrt(v) * 100;                   % std dev in percent
corr = vcv ./ sqrt(v' * v);             % correlation matrix

corrWithY = corr(:, 1);                 % correlation with output (column 1)

% persistence of output: autocorrelation
y1 = dev(2:Tq, 1);                      % periods 2:T
y0 = dev(1:Tq-1, 1);                    % periods 1:T-1
ac = corrcoef(y1, y0);
rho = ac(1, 2);
\end{lstlisting}

\lstinline[style=inl]!dev'*dev! multiplies the $202\times14$ matrix by its transpose, giving
$14\times14$. \lstinline[style=inl]!sum(vcv .* eye(14))! keeps only the diagonal (the variances);
\lstinline[style=inl]!sqrt(v'*v)! is the outer product of the standard deviations, which is
exactly the denominator of every correlation.

\section{Plotting}

\begin{lstlisting}
t = (1:Tq)';
plot(t, dev(:,1), "LineWidth", 1.2);  hold on
plot(t, dev(:,2), "LineWidth", 1.2);  hold off
xticks(t(1:16:end))                     % one tick every 4 years
xticklabels(q(1:16:end))                % ... labelled with the quarter strings
xlabel("quarter"); ylabel("log deviation from trend")
legend(["Output", "Consumption"], Location = "best"); grid on

figure
yyaxis left;  plot(t, trend(:,1));  ylabel("trend")
yyaxis right; plot(t, dev(:,1));    ylabel("cycle")

exportgraphics(gcf, "fig.png", Resolution = 200)   % or saveas(gcf,"fig.png")
\end{lstlisting}

\lstinline[style=inl]!xticks! needs \emph{numbers}, so pass positions (\lstinline[style=inl]!t(1:16:end)!) and give
the labels separately --- calling \lstinline[style=inl]!xticks! on the quarter strings errors.

\section{Function quick reference}

{\footnotesize
\begin{center}
\begin{tabular}{@{}>{\raggedright\arraybackslash}p{4.1cm}>{\raggedright\arraybackslash}p{6.3cm}>{\raggedright\arraybackslash}p{6.3cm}@{}}
\toprule
task & syntax & note \\
\midrule
make a vector & \texttt{zeros(Tq,1)}, \texttt{ones(Tq,1)}, \texttt{(1:Tq)'} & \texttt{'} transposes \\
size & \texttt{size(X)}, \texttt{size(X,1)}, \texttt{height(T)} & \texttt{numel} for total count \\
mean / sum & \texttt{mean(x)}, \texttt{sum(x)}, \texttt{mean(M)} & on a matrix: column-wise \\
std & \texttt{std(x,1)} & \texttt{1} = divide by $N$ (population) \\
growth & \texttt{diff(x)./x(1:end-1)} & \texttt{diff} is one shorter \\
cumulate & \texttt{cumsum(x)} & \\
logs & \texttt{log(x)}, \texttt{exp(x)} & natural log \\
correlation & \texttt{corrcoef(a,b)}, \texttt{corr(M)} & $2\times2$ / matrix \\
diagonal, identity & \texttt{diag(v)}, \texttt{eye(14)} & \\
inverse, solve & \texttt{inv(A)}, \texttt{A\textbackslash{}b} & \texttt{A\textbackslash{}b} is faster \\
indexing & \texttt{x(1)}, \texttt{X(:,k)}, \texttt{X(1:5,2)}, \texttt{x(end)} & starts at 1, not 0 \\
text & \texttt{string(x)}, \texttt{strtrim(s)}, \texttt{join(v,"\_")} & \texttt{==} works on strings \\
search in names & \texttt{contains(var,"goods")}, \texttt{startsWith}, \texttt{find(...)} & useful to find a column \\
print & \texttt{fprintf("\%d rows\textbackslash{}n", Tq)} & \texttt{\%g}, \texttt{\%.2f} \\
write files & \texttt{writematrix(X,"x.csv")}, \texttt{writetable} & \\
sanity check & \texttt{assert(numel(y)==Tq)}, \texttt{isnan(sum(y))} & catch size errors early \\
clear the deck & \texttt{clear}, \texttt{clc}, \texttt{close all} & start of every script \\
\bottomrule
\end{tabular}
\end{center}}

\section{Errors I hit, and what fixes them}

{\footnotesize
\begin{center}
\begin{tabular}{@{}>{\raggedright\arraybackslash}p{6.6cm}>{\raggedright\arraybackslash}p{9.9cm}@{}}
\toprule
message / symptom & fix \\
\midrule
\emph{Invalid text character} & the name you typed contains something MATLAB cannot use (a space, a dash, a \texttt{\$}). Use \texttt{T.("...")} or fix the header \\
\emph{VariableNames must be a character vector} & use name=value pairs: \texttt{table(a, VariableNames=["x"])}, not \texttt{table(a,"VariableNames",["x"])} \\
\emph{Incorrect dimensions for matrix multiplication} & you meant \texttt{.*} or \texttt{./} instead of \texttt{*} / \texttt{/} \\
\texttt{xticks} error on a categorical axis & pass numeric positions, then \texttt{xticklabels(q(...))} \\
\texttt{height} fails on a struct & use \texttt{numel()} or keep the table instead of extracting fields \\
results look absurd (huge trend swings) & you skipped \texttt{log()} before filtering \\
\texttt{NaN} in the first row of a growth series & expected: \texttt{diff} shortens the series; index with \texttt{(1:end-1)} \\
numbers off by one quarter & MATLAB indexes from 1, an "0-based" habit from other languages \\
script name starts with \texttt{\_} or a digit & MATLAB scripts need a valid identifier as filename \\
\bottomrule
\end{tabular}
\end{center}}

\section{The whole pipeline, in order}

\begin{lstlisting}
load and inspect  ->  build the 14 series  ->  per capita  ->  log
   ->  HP filter (trend + cycle)  ->  variance-covariance  ->  std dev
   ->  correlations with output  ->  autocorrelation of output
   ->  plot trend & cycle  ->  write the table / figure out
\end{lstlisting}

Keep the steps in separate cells (or \lstinline[style=inl]!%%! sections) so you can rerun from
the middle without recomputing everything.

\end{document}
